\begin{abstract}
We examine whether enterprise identity loss can be formalized as a phase
transition rather than as organizational failure for enterprises modeled as
continua $K_{\mathrm{EA}}$ within the Ontology of Continua and its executable
specification ETS.K$_{\mathrm{EA}}$. Enterprise identity is defined as a strictly
structural property: an admissible identity-supporting subspace sustained by a
non-empty set of invariant cycles. Identity loss occurs when this supporting
structure loses admissibility or when its sustaining cycles are extinguished,
while the enterprise continuum itself remains admissible and operational.
Accordingly, identity loss constitutes a distinct phase transition characterized
by a discontinuous change in invariant structure rather than by collapse of the
admissible state space. Conditions for irreversibility, formal no-go results for
identity reconstitution under closed dynamics, and explicit falsifiability
criteria are derived. The analysis is purely structural and introduces no new
primitives, normative claims, or managerial interpretations.
\end{abstract}

\noindent\textbf{Keywords:} enterprise identity; phase transition; Ontology of
Continua; admissibility; invariant cycles; irreversibility
